\section{Separate parameters only when the experiment separates directions}
\label{sec:identification}
The following is a conditional model analysis. Neither experimental file
is used here to estimate a rotor-angle error or to jointly identify physical
resistance and inverter parameters. The failed global constant-resistance
description motivates stating which additional information would make such
a separation meaningful.
At one pair, changing resistance moves the residual vector along
$\vect s_R$; changing angle moves it along $\vect s_\gamma$. If these
directions are parallel, the observation cannot tell which error caused the
motion. Collecting many pairs can remove this ambiguity because the two
directions change differently with current and speed.

\subsection{Stacking, rank, and conditioning}
Let $\vect\theta=[\resistanceerror,\angleerror]\trans$ and
$\matr S_k=[\vect s_{R,k}\ \vect s_{\gamma,k}]$. Then
\begin{equation}
 \vect r_k\approx\matr S_k\vect\theta+\vect\eta_k,\qquad
 \vect r=\begin{bmatrix}\vect r_1\\\vdots\\\vect r_N\end{bmatrix},\quad
 \matr S=\begin{bmatrix}\matr S_1\\\vdots\\\matr S_N\end{bmatrix},
 \quad\vect r\approx\matr S\vect\theta+\vect\eta.
 \label{eq:stacked}
\end{equation}
Here $\vect\eta$ includes noise, interpolation error, neglected sources, and
linearization discrepancy. Distinct small parameter changes are distinguishable
in this conditional two-parameter model exactly when
\begin{equation}
 \boxed{\rank(\matr S)=2.}\label{eq:rank-condition}
\end{equation}
This is local identifiability given the magnetic map and nuisance assumptions;
it is not a guarantee of globally unique calibration. In particular, angle
sensitivities contain unknown ideal flux and derivatives. They must come from
a calibrated magnetic reference or a sufficiently constrained joint fit.
Computing derivatives from an unconstrained erroneous map can bias the columns.
Symmetric projections provide an initialization for small offsets, but do not
solve the full calibration problem by themselves.

Near-parallel columns make the separation sensitive to noise even with rank
two. Before interpreting a singular-value condition number, whiten residuals
by their covariance and scale parameters by meaningful error magnitudes:
\begin{equation}
 \matr A=\matr C^{-1/2}\matr S\matr D,\qquad
 \matr D=\diag(R_{\rm scale},\gamma_{\rm scale}).
\end{equation}
Singular values of $\matr A$ measure distinguishability of scaled parameter
changes in noise units \cite{horn2013}. A condition number of unscaled columns
mixing ohms and radians has no invariant engineering meaning. Pair covariance
must reflect common samples and any shared interpolation.

Only after this check, weighted least squares is appropriate:
\begin{equation}
 \hat{\vect\theta}=\arg\min_{\vect\theta}
 (\vect r-\matr S\vect\theta)\trans\matr C^{-1}
 (\vect r-\matr S\vect\theta).\label{eq:least-squares}
\end{equation}
Use QR or SVD numerically. If the linear model and covariance are valid,
$(\matr S\trans\matr C^{-1}\matr S)^{-1}$ describes parameter covariance;
model bias is not covered by that expression. Robust loss functions can limit
outlier influence but cannot restore absent rank.

\subsection{An isotropic thought experiment and its nearby cases}
At $i_d=0$ in a linear isotropic PM machine,
\begin{equation}
 r_d=-\frac{i_q}{\omegae}\resistanceerror,\qquad
 r_q=-\psi_{\rm pm}\angleerror.\label{eq:isotropic}
\end{equation}
The channels decouple. A pair with $i_q\ne0$, $\omegae\ne0$, and
$\psi_{\rm pm}\ne0$ has rank two. Set all three quantities to one in
consistent normalized units: $[r_d,r_q]=[-\resistanceerror,-\angleerror]$. This is the
cleanest picture of what each channel measures.

Add saliency: $r_d$ now also contains $(L_q-L_d)i_q\angleerror$, while $r_q$
retains its PM angle signal at $i_d=0$. If only d flux is retained at one speed,
resistance and angle are indistinguishable in this linear sweep. Add d/q
saturation: local gains and q flux change the angle column. Add cross-saturation:
the mixed gains further couple its dependence on both currents. The isotropic
separation becomes less simple, but \cref{eq:columns} still applies.

If PM flux vanishes in an isotropic linear machine, a rotation leaves its
entire map unchanged and $\vect s_\gamma=0$. No experiment based on this map
can identify its orientation. At $i_q=i_d=0$ the resistance column is zero.
Collecting many such uninformative points cannot repair these degeneracies.

\subsection{Speed provides an additional experimental direction}
At repeated current labels, constant thermal state, and quasi-static
magnetics, the resistance column scales with $1/\omegae$, whereas a constant
geometric angle column does not. Thus $\omegae r_d$ is speed-independent for
a pure constant resistance error. If one scalar current-pair channel has
nonzero resistance and angle sensitivities, two distinct speeds can separate
them because the two rows have different resistance scaling.
\begin{figure}[tb]
\centering
\begin{tikzpicture}
\begin{axis}[diagnostic axis,xlabel={$\normalizedspeed=\omega_e/\omega_*$},
 title={Fixed current, separate error injections},ylabel={$r/\psi_*$},
 xmin=.5,xmax=2,ymin=-.045,ymax=.012]
 \addplot[estimated quantity,strong line] table[x=speed,y=rdR]{figures/data/speeds.dat};
 \addlegendentry{$r_{d,R}$}
 \addplot[torque quantity,standard line] table[x=speed,y=rdg]{figures/data/speeds.dat};
 \addlegendentry{$r_{d,\gamma}$}
 \addplot[reference quantity,alternative pattern,standard line] table[x=speed,y=rqg]{figures/data/speeds.dat};
 \addlegendentry{$r_{q,\gamma}$}
\end{axis}
\begin{axis}[diagnostic axis,at={(.5\textwidth,0)},anchor=south west,
 xlabel={$\normalizedspeed=\omega_e/\omega_*$},title={Remove the inverse-speed factor},
 ylabel={$\normalizedspeed r_{d,R}/\psi_*$},xmin=.5,xmax=2,ymin=-.025,ymax=-.015]
 \addplot[estimated quantity,strong line] table[x=speed,y=scaledR]{figures/data/speeds.dat};
\end{axis}
\end{tikzpicture}
\caption{Exact synthetic residuals at $(x,y)=(0,0.8)$ for separate injections
$\resistanceerror/R_*=0.025$ and $\angleerror=0.02$ rad. Resistance scales inversely with
speed; constant geometric misalignment does not. The right plot is a scaling
check, not proof of resistance rather than inverter error.}
\label{fig:speed}
\end{figure}


This inference requires the same magnetic and error parameters across those
speeds. Inverter voltage errors also enter through $1/\omegae$; a delay can
produce $\angleerror=\gamma_0+\omegae\timingdelay$ under a convention in which positive
$\timingdelay$ is a phase advance. A physical sensing lag has the opposite sign.
If both are unknown, include a third column $\omegae\vect s_\gamma$ and
test rank three. Frequency-dependent effective resistance, magnetic losses,
or changing temperatures invalidate the simple two-column comparison.

An unknown current-proportional voltage error is an exact counterexample to
unique resistance identification:
\begin{equation}
 \vect\varepsilon_u=-a\current
 \quad\Longrightarrow\quad
 \delta\flux_u=\frac{a}{\omegae}\rotationgenerator\current.
 \label{eq:confounding}
\end{equation}
It has precisely the fingerprint of $\resistanceerror=a$, at every speed and current.
If nuisance columns form $\matr B$, separate identification requires
$\rank([\matr B\ \matr S])-\rank(\matr B)=2$. Extra speeds alone cannot
remove a nuisance that occupies the same column space. Position-offset
injection is a related experimental parameter-identification technique
\cite{liu2017}; the present method analyzes symmetry residuals rather than
claiming to reproduce that estimator.
